% theme: strata_and_sedimentary_rocks
\MajorQuestion{地層と堆積岩}
\SetRowSpaceQanda

\Instruction{次の問いに答えなさい。}
\QQ{運ばれた土砂が水底などに積もることを\NamedBlank{deposition}という。}{堆積}
\QQ{水底などに積もった物質が、上からの重みで押し固められてできた岩石を\NamedBlank{sedimentary_rock}という。}{堆積岩}
\QQ{生物の死がいや生活の跡が、重なった層の中に残ったものを\NamedBlank{fossil}という。}{化石}
\QQ{岩石が温度変化や水などの影響を受け、表面からしだいにくずれていくことを何というか。}{風化}
\QQ{ある地点の層の上下関係や厚さ、特徴を柱のように表した図を何というか。}{柱状図}

\Instruction{\strut}
\QQ{川によって運ばれた土砂が河口付近に\RefBlank{deposition}してできる、低く平らな地形を何というか。}{三角州}
\QQ{生物の死がいなどが積もってでき、うすい塩酸をかけると二酸化炭素が発生する\RefBlank{sedimentary_rock}を何というか。}{石灰岩}
\QQ{流水が岩石や土地をけずるはたらきを何というか。}{侵食}
\QQ{離れた場所の層を比べるときに目印となる、特徴的な層を何というか。}{鍵層}
\QQ{\RefBlank{fossil}のうち、層ができた当時の環境を知る手がかりとなるものを何というか。}{示相化石}

\Instruction{\strut}
\QQ{直径$2$\,mm以上の粒が押し固められてできた\RefBlank{sedimentary_rock}を何というか。}{れき岩}
\QQ{\RefBlank{sedimentary_rock}のうち、生物の死がいなどが積もってでき、うすい塩酸をかけても気体が発生しない硬いものを何というか。}{チャート}
\QQ{流水がけずり取った土砂を運ぶはたらきを何というか。}{運搬}
\QQ{横から強い力を受けて、波を打つように曲がった層のつくりを何というか。}{しゅう曲}
\QQ{\RefBlank{fossil}のうち、層ができた年代を知る手がかりとなるものを何というか。}{示準化石}

\Instruction{\strut}
\QQ{直径$\dfrac{1}{16}$\,mm以上$2$\,mm未満の粒が押し固められてできた\RefBlank{sedimentary_rock}を何というか。}{砂岩}
\QQ{\RefBlank{sedimentary_rock}のうち、火山灰などの火山噴出物が積もってできたものを何というか。}{凝灰岩}
\QQ{重なった層が地表に現れているところを何というか。}{露頭}
\QQ{大地にはたらく強い力によって、層が切れてずれたものを何というか。}{断層}
\QQ{おもに直径$\dfrac{1}{16}$\,mm未満の粒が押し固められてできた\RefBlank{sedimentary_rock}を何というか。}{泥岩}
